% CPSY 1291 -- Linear algebra for this course
% A handout to help. Written for a student who found the linear algebra in the first
% lectures fast. It defines every symbol the lectures and Assignment 1 use, with a small
% numerical example per concept. There is no code here beyond a one-line NumPy aside.
% Compile:  latexmk -lualatex linear-algebra-notes.tex
\documentclass[11pt]{article}
\usepackage{cpsy1291-handout}

\newcommand{\vx}{\bm{x}}
\newcommand{\vy}{\bm{y}}
\newcommand{\vw}{\bm{w}}
\newcommand{\mX}{\bm{X}}
\newcommand{\mD}{\bm{D}}
\newcommand{\norm}[1]{\left\lVert #1 \right\rVert}
\newcommand{\deuc}{d_{\mathrm{euc}}}
\newcommand{\dcos}{d_{\cos}}
\newcommand{\dcorr}{d_{\mathrm{corr}}}
\newcommand{\np}[1]{{\small\color{lightgray}\textit{In NumPy:} \texttt{#1}}\par}

\begin{document}

\begin{center}
{\sffamily\small\color{accent}\MakeUppercase{CPSY 1291 \ $\cdot$\ A handout to help}}\\[5pt]
{\sffamily\LARGE\bfseries\color{brown} Linear algebra for this course}\\[3pt]
{\sffamily\large\color{ink} Vectors, distances, and dissimilarity matrices, with the numbers written out}\\[6pt]
{\color{lightgray}\small Fall 2026 \quad$\cdot$\quad read at your own pace; every symbol here is the one the lectures use}
\end{center}
\vspace{4pt}\hrule\vspace{10pt}

\section*{Who this is for}
\addcontentsline{toc}{section}{Who this is for}

Several of you said that the linear algebra in the first lectures went by quickly.
This handout slows it down. It contains nothing that is not already in the lectures;
it just takes each object one at a time, defines it, and works a small numerical example
by hand so that you can see what the symbols are doing. Vectors here have two or three
entries so the arithmetic fits on a line. The vectors in the course have thousands of
entries, and nothing changes except the length of the sums.

\begin{keybox}
\ask{The one idea.} A list of $D$ numbers is a point in a $D$-dimensional space. Once
two things are points, ``how similar are they'' becomes ``how far apart are they'', and
there is more than one reasonable answer to that. Everything below is about making that
sentence precise.
\end{keybox}

\section{Vectors, and the names the course uses}

\subsection{A vector is a list of numbers, and also a point}

A \emph{vector} is an ordered list of numbers. Write it in bold, $\vx$, and write its
$j$-th entry as $x_j$:
\[
\vx = (x_1, x_2, \dots, x_D) .
\]
$D$ is the number of entries, the \emph{dimension}. Two readings of the same list are
worth holding at once. As a \textbf{point}, $\vx$ is a location: with $D=2$, $(3, 4)$ is
the point three units right and four units up. As an \textbf{arrow}, $\vx$ is the
displacement from the origin to that point, with a length and a direction. The lectures use both.

\textbf{Example.} Record the firing rates of three neurons to one image and get
$\vx = (2, 0, 5)$: neuron 1 fired at 2 spikes per second, neuron 2 was silent, neuron 3
fired at 5. That image is now a point in a three-dimensional ``neuron space'', one axis per
neuron. Pixel values, a row of network activations, and a set of ratings are all vectors
in exactly this sense.

\subsection{Stimuli are indexed with a superscript in parentheses}

The course looks at many stimuli at once, so it needs two indices: \emph{which stimulus}
and \emph{which entry}. The convention, used in every lecture and assignment:
\begin{itemize}
\item $\vx^{(i)}$ is the response vector to stimulus $i$, for $i = 1, \dots, N$;
\item $x^{(i)}_j$ is its $j$-th entry, for $j = 1, \dots, D$.
\end{itemize}
So $N$ counts stimuli (images, trials, conditions) and $D$ counts entries (pixels, units,
neurons, rating scales). Read $x^{(7)}_{3}$ as ``the response of neuron 3 to image 7''.

\subsection{Stacking the vectors: the response matrix}

Put the $N$ response vectors on top of each other, one per row, and you get the
\emph{response matrix} $\mX$, with $N$ rows and $D$ columns:
\[
\mX =
\begin{pmatrix}
\text{---}\ \vx^{(1)}\ \text{---}\\
\text{---}\ \vx^{(2)}\ \text{---}\\
\vdots\\
\text{---}\ \vx^{(N)}\ \text{---}
\end{pmatrix}
\qquad
\mX \text{ is } N \times D, \qquad
(\mX)_{ij} = x^{(i)}_j .
\]
\textbf{Example.} Three images, two neurons ($N = 3$, $D = 2$), with
$\vx^{(1)} = (3, 4)$, $\vx^{(2)} = (4, 3)$, $\vx^{(3)} = (0, 1)$:
\[
\mX = \begin{pmatrix} 3 & 4 \\ 4 & 3 \\ 0 & 1 \end{pmatrix}, \qquad
\text{row 2 is } \vx^{(2)},\quad \text{column 2 is neuron 2's response to every image.}
\]
``$N \times D$'' always means rows first, then columns. \hl{Rows are stimuli; columns are
features.} Every data matrix in this course is laid out this way.

\np{X.shape is (N, D); X[i] is row i; X[:, j] is column j.}

\section{Adding, subtracting, and scaling vectors}

Vectors of the same dimension are added and subtracted entry by entry, and multiplied by
a number (a \emph{scalar}) entry by entry:
\[
\vx + \vy = (x_1 + y_1, \dots, x_D + y_D), \qquad
\vx - \vy = (x_1 - y_1, \dots, x_D - y_D), \qquad
c\,\vx = (c x_1, \dots, c x_D).
\]
\textbf{Example.} With $\vx = (3, 4)$ and $\vy = (4, 3)$:
\[
\vx + \vy = (7, 7), \qquad \vx - \vy = (-1, 1), \qquad 2\vx = (6, 8), \qquad
\tfrac{1}{2}\vx = (1.5, 2).
\]
Geometrically, $\vx + \vy$ is the arrow you get by placing $\vy$'s tail at $\vx$'s head;
$c\vx$ points the same way as $\vx$ (the opposite way if $c < 0$) and is $|c|$ times as
long.

\begin{keybox}
\ask{The difference vector.} $\vx - \vy$ is the arrow \emph{from} the point $\vy$ \emph{to}
the point $\vx$. Its length is the distance between them. That one sentence is the whole
of \S4, so it is worth a second read: subtracting two response vectors gives the
displacement between two stimuli in neuron space.
\end{keybox}

\np{x + y, x - y, 2 * x, all entry by entry.}

\section{The dot product}

\subsection{Definition}

The \emph{dot product} of two vectors of the same dimension multiplies matching entries
and adds the results. It turns two vectors into one number:
\[
\vx \cdot \vy = \sum_{j=1}^{D} x_j\,y_j = x_1 y_1 + x_2 y_2 + \dots + x_D y_D .
\]
\textbf{Example.} $(3, 4)\cdot(4, 3) = 3\cdot4 + 4\cdot3 = 12 + 12 = 24$.\quad
$(3, 4)\cdot(0, 1) = 0 + 4 = 4$.\quad $(3, 4)\cdot(4, -3) = 12 - 12 = 0$.

The dot product is symmetric, $\vx \cdot \vy = \vy \cdot \vx$, and a vector dotted with
itself gives the sum of its squared entries: $(3, 4)\cdot(3, 4) = 9 + 16 = 25$.

\subsection{The Euclidean norm, or length}

The \emph{norm} of $\vx$ is its length as an arrow, by Pythagoras:
\[
\norm{\vx} = \sqrt{\vx \cdot \vx} = \sqrt{\textstyle\sum_j x_j^2} .
\]
\textbf{Example.} $\norm{(3, 4)} = \sqrt{9 + 16} = 5$.\quad
$\norm{(4, 3)} = 5$.\quad $\norm{(0, 1)} = 1$.\quad $\norm{(1, 1, 1)} = \sqrt{3} \approx 1.732$.

Scaling a vector scales its length: $\norm{c\vx} = |c|\,\norm{\vx}$, so
$\norm{2\,(3, 4)} = \norm{(6, 8)} = 10$. Dividing a vector by its own norm gives a
\emph{unit vector}, length exactly 1, pointing the same way: $(3, 4)/5 = (0.6, 0.8)$,
and indeed $0.36 + 0.64 = 1$.

\np{np.linalg.norm(x); np.dot(x, y) or x @ y for the dot product.}

\subsection{The geometric meaning: the angle between two vectors}

The dot product is also determined by the two lengths and the angle $\phi$ between the
arrows:
\[
\vx \cdot \vy = \norm{\vx}\,\norm{\vy}\,\cos\phi ,
\qquad\text{equivalently}\qquad
\cos\phi = \frac{\vx \cdot \vy}{\norm{\vx}\,\norm{\vy}} .
\]
This is the fact that makes the dot product useful. Read it as: \emph{how much of $\vx$
points along $\vy$}, scaled by both lengths. It is large and positive when the two arrows
point the same way ($\phi$ near $0$, $\cos\phi$ near $1$), zero when they are
perpendicular ($\phi = 90^\circ$, $\cos\phi = 0$), and negative when they point in
opposite directions ($\phi$ near $180^\circ$, $\cos\phi$ near $-1$).

\textbf{Example.} For $\vx = (3, 4)$ and $\vy = (4, 3)$,
\[
\cos\phi = \frac{24}{5 \cdot 5} = 0.96, \qquad \phi = \arccos(0.96) \approx 16.3^\circ .
\]
For $\vx = (3, 4)$ and $\vy = (4, -3)$ the dot product was $0$, so $\cos\phi = 0$ and the
arrows are perpendicular; a sketch confirms it. For $\vx = (3, 4)$ and $2\vx = (6, 8)$,
$\cos\phi = 50 / (5 \cdot 10) = 1$: same direction, angle zero.

\section{Euclidean distance}

The \emph{Euclidean distance} between two points is the length of the difference vector,
the straight-line distance between them:
\[
\deuc(\vx^{(i)}, \vx^{(k)}) = \norm{\vx^{(i)} - \vx^{(k)}}
= \sqrt{\textstyle\sum_j \big(x^{(i)}_j - x^{(k)}_j\big)^2} .
\]
\textbf{Example.} With $\vx^{(1)} = (3, 4)$, $\vx^{(2)} = (4, 3)$, $\vx^{(3)} = (0, 1)$:
\begin{align*}
\deuc(\vx^{(1)}, \vx^{(2)}) &= \norm{(-1, 1)} = \sqrt{2} \approx 1.414, \\
\deuc(\vx^{(1)}, \vx^{(3)}) &= \norm{(3, 3)} = \sqrt{18} \approx 4.243, \\
\deuc(\vx^{(2)}, \vx^{(3)}) &= \norm{(4, 2)} = \sqrt{20} \approx 4.472 .
\end{align*}
So images 1 and 2 are close, and image 3 is far from both. This is the distance you would
measure with a ruler on a drawing of the three points.

Three properties, all visible in the example: the distance from a point to itself is
$0$; $\deuc(\vx, \vy) = \deuc(\vy, \vx)$; and a distance is never negative. They matter
in \S7.

Euclidean distance is sensitive to overall magnitude. A neuron population that fires twice
as hard to every image moves every point twice as far from the origin and doubles every
distance, even though the \emph{pattern} across neurons is unchanged. The next two
measures are designed to ignore that.

\np{np.linalg.norm(x - y).}

\section{Cosine similarity and cosine dissimilarity}

\emph{Cosine similarity} is just $\cos\phi$ from \S3.3: the dot product with both lengths
divided out. It compares directions and ignores lengths entirely. The course turns it
into a dissimilarity (small means similar) by subtracting from one:
\[
\dcos(\vx^{(i)}, \vx^{(k)}) = 1 - \cos\phi
= 1 - \frac{\vx^{(i)} \cdot \vx^{(k)}}{\norm{\vx^{(i)}}\,\norm{\vx^{(k)}}} .
\]
It is $0$ when the two vectors point the same way, $1$ when they are perpendicular, and
$2$ when they point in opposite directions; so it lives in $[0, 2]$.

\textbf{Example.} $\vx^{(1)} = (3, 4)$ and $\vx^{(2)} = (4, 3)$: $\cos\phi = 0.96$, so
$\dcos = 0.04$. Nearly the same direction; nearly zero dissimilarity.

\textbf{Example, where the two measures disagree.} Take $\vx = (1, 0)$ and $\vy = (3, 0)$.
Euclidean distance: $\norm{(-2, 0)} = 2$. Cosine: $\cos\phi = 3 / (1 \cdot 3) = 1$, so
$\dcos = 0$. Euclidean says ``two units apart''; cosine says ``identical''. Neither is
wrong: one asks \emph{how far apart are the points}, the other \emph{do the arrows point
the same way}. If $\vy$ is the same neurons firing three times as hard, cosine is the
measure that says so.

\np{1 - np.dot(x, y) / (np.linalg.norm(x) * np.linalg.norm(y)), for one pair of vectors.}

\section{Correlation and correlation dissimilarity}

\subsection{Centring a vector by its own mean}

The \emph{mean} of a vector is the average of its entries, $\bar{x} = \frac{1}{D}\sum_j x_j$.
\emph{Centring} subtracts it from every entry, so the centred vector
$\tilde{\vx} = \vx - \bar{x}$ has mean zero. This is a per-vector operation: each stimulus
is centred by its \emph{own} mean, across its $D$ entries.

\textbf{Example.} $\vx = (1, 2, 3)$ has mean $2$, so $\tilde{\vx} = (-1, 0, 1)$.
$\vy = (2, 4, 6)$ has mean $4$, so $\tilde{\vy} = (-2, 0, 2)$.
$\bm{z} = (3, 2, 1)$ has mean $2$, so $\tilde{\bm{z}} = (1, 0, -1)$.

\subsection{Pearson correlation is cosine similarity after centring}

The \emph{Pearson correlation} $r$ between two vectors is the cosine of the angle between
their centred versions:
\[
r(\vx, \vy) = \frac{\tilde{\vx} \cdot \tilde{\vy}}{\norm{\tilde{\vx}}\,\norm{\tilde{\vy}}}
= \frac{\sum_j (x_j - \bar{x})(y_j - \bar{y})}
       {\sqrt{\sum_j (x_j - \bar{x})^2}\,\sqrt{\sum_j (y_j - \bar{y})^2}} .
\]
The second form is the textbook definition; the first says what it is: cosine on centred
vectors. $r$ lies in $[-1, 1]$. It is $1$ when one vector is a positive-slope straight-line
function of the other, $-1$ for a negative slope, and $0$ when there is no linear relation.

\textbf{Example.} $\tilde{\vx} = (-1, 0, 1)$ and $\tilde{\vy} = (-2, 0, 2)$:
$r = (2 + 0 + 2) / (\sqrt{2}\cdot\sqrt{8}) = 4 / 4 = 1$. Indeed $\vy = 2\vx$.
$\tilde{\vx} = (-1, 0, 1)$ and $\tilde{\bm{z}} = (1, 0, -1)$:
$r = (-1 + 0 - 1) / (\sqrt{2}\cdot\sqrt{2}) = -1$. The pattern in $\bm{z}$ is the reverse
of the pattern in $\vx$.

The course's third dissimilarity is
\[
\dcorr(\vx^{(i)}, \vx^{(k)}) = 1 - r ,
\]
again in $[0, 2]$: $0$ for a perfectly matching pattern, $1$ for no linear relation, $2$
for a perfectly reversed one.

\begin{keybox}
\ask{Why centre?} Cosine ignores how \emph{strongly} a population fires overall.
Correlation also ignores its \emph{baseline}: adding the same constant to every neuron's
rate leaves $r$ unchanged, because centring removes it. Correlation therefore compares
only the \emph{pattern} of ups and downs across neurons.
\end{keybox}

\textbf{Example, where cosine and correlation disagree.} Take $\vx = (5, 5, 6)$ and
$\vy = (6, 5, 5)$. Both sit far from the origin in almost the same direction:
$\vx \cdot \vy = 30 + 25 + 30 = 85$, $\norm{\vx} = \norm{\vy} = \sqrt{86}$, so
$\cos\phi = 85/86 \approx 0.988$ and $\dcos \approx 0.012$. Now centre. Both means are
$16/3$, so
$\tilde{\vx} = (-\tfrac13, -\tfrac13, \tfrac23)$ and
$\tilde{\vy} = (\tfrac23, -\tfrac13, -\tfrac13)$. Then
$\tilde{\vx} \cdot \tilde{\vy} = -\tfrac29 + \tfrac19 - \tfrac29 = -\tfrac13$, each centred
norm is $\sqrt{6/9}$, and
\[
r = \frac{-1/3}{6/9} = -0.5, \qquad \dcorr = 1.5 .
\]
Cosine calls them nearly identical because both are dominated by a large shared baseline
(every neuron firing at about 5). Correlation strips the baseline and sees what is left:
$\vx$ peaks on neuron 3, $\vy$ peaks on neuron 1, opposite patterns. When the vectors are
ReLU activations, which are non-negative and often share a large common offset, this
difference between the two measures is not a technicality.

\np{np.corrcoef(x, y)[0, 1] gives r for one pair of vectors.}

\subsection{The three dissimilarities, side by side}

For two response vectors $\vx^{(i)}$ and $\vx^{(k)}$:
\[
\deuc = \norm{\vx^{(i)} - \vx^{(k)}}, \qquad
\dcos = 1 - \cos\phi, \qquad
\dcorr = 1 - r .
\]
\begin{center}
\small
\begin{tabular}{@{}lllll@{}}
\toprule
measure & ignores & range & $0$ means & example, $(3,4)$ vs.\ $(4,3)$ \\
\midrule
$\deuc$ & nothing & $[0, \infty)$ & same point & $\sqrt{2} \approx 1.414$ \\
$\dcos$ & overall magnitude & $[0, 2]$ & same direction & $0.04$ \\
$\dcorr$ & magnitude and baseline & $[0, 2]$ & same pattern & $2$ \\
\bottomrule
\end{tabular}
\end{center}
The last cell is worth checking by hand: $(3, 4)$ centres to $(-0.5, 0.5)$ and $(4, 3)$ to
$(0.5, -0.5)$, exactly opposite, so $r = -1$ and $\dcorr = 2$. With only two entries a
pattern can only go up or down, so $D = 2$ correlations are always $\pm 1$; with
$D = 4096$ the measure is far more graded.

\section{Matrices}

\subsection{What a matrix is, and two ways to read one}

A matrix is a rectangular grid of numbers with $m$ rows and $n$ columns, ``$m \times n$''.
Its entry in row $i$, column $j$ is written with two subscripts, row first: $(\bm{A})_{ij}$
or $a_{ij}$. Two readings recur in this course.
\textbf{As data:} $\mX$ from \S1.3, one row per stimulus, one column per feature.
\textbf{As a map:} a matrix $\bm{W}$ turns a vector into a new vector, $\vw \mapsto
\bm{W}\vw$, which is what a layer of a neural network does.

\subsection{The matrix--vector product is a stack of dot products}

If $\bm{A}$ is $m \times n$ and $\vw$ has $n$ entries, the product $\bm{A}\vw$ is a vector
with $m$ entries, and \hl{its $i$-th entry is the dot product of row $i$ of $\bm{A}$ with
$\vw$}:
\[
(\bm{A}\vw)_i = \sum_{j=1}^{n} a_{ij}\,w_j = \bm{a}_i \cdot \vw ,
\quad\text{where } \bm{a}_i \text{ is row } i .
\]
\textbf{Example.} With the $3 \times 2$ response matrix from \S1.3 and $\vw = (1, -1)$:
\[
\mX\vw = \begin{pmatrix} 3 & 4 \\ 4 & 3 \\ 0 & 1 \end{pmatrix}
\begin{pmatrix} 1 \\ -1 \end{pmatrix}
= \begin{pmatrix} 3 - 4 \\ 4 - 3 \\ 0 - 1 \end{pmatrix}
= \begin{pmatrix} -1 \\ 1 \\ -1 \end{pmatrix} .
\]
Read it as ``project every stimulus onto the direction $\vw$'': one number per image.
The inner dimensions must agree: a $3 \times 2$ matrix takes a vector with $2$ entries and
returns one with $3$. A matrix times a matrix works the same way, one column at a time:
$(m \times k)\,(k \times n) \to (m \times n)$.

\np{A @ w; the shapes must read (m, n) @ (n,) -> (m,).}

\subsection{The transpose}

The \emph{transpose} $\bm{A}^{\top}$ swaps rows and columns: $(\bm{A}^{\top})_{ij} = a_{ji}$,
and an $m \times n$ matrix becomes $n \times m$.
\[
\mX = \begin{pmatrix} 3 & 4 \\ 4 & 3 \\ 0 & 1 \end{pmatrix}
\quad\Longrightarrow\quad
\mX^{\top} = \begin{pmatrix} 3 & 4 & 0 \\ 4 & 3 & 1 \end{pmatrix} .
\]
Transposing a data matrix turns ``one row per stimulus'' into ``one column per stimulus''.
One product is worth knowing by name:
$\mX\mX^{\top}$ is $N \times N$ and its entry $(i, k)$ is the dot product of row $i$ with
row $k$, that is, $\vx^{(i)} \cdot \vx^{(k)}$, every pairwise dot product at once.
\[
\mX\mX^{\top} =
\begin{pmatrix} 3 & 4 \\ 4 & 3 \\ 0 & 1 \end{pmatrix}
\begin{pmatrix} 3 & 4 & 0 \\ 4 & 3 & 1 \end{pmatrix}
=
\begin{pmatrix} 25 & 24 & 4 \\ 24 & 25 & 3 \\ 4 & 3 & 1 \end{pmatrix} .
\]
The diagonal holds the squared norms ($25, 25, 1$, matching \S3.2) and the off-diagonal
entries are the dot products computed in \S3.1. A matrix that equals its own transpose,
as this one does, is called \emph{symmetric}.

\np{X.T; X @ X.T.}

\section{The dissimilarity matrix}

\subsection{Definition}

Choose one of the three measures and compute it for \emph{every pair} of stimuli. The
result is a square $N \times N$ matrix $\mD$, the \emph{dissimilarity matrix} (the
lectures also call it the RDM, representational dissimilarity matrix):
\[
d_{ik} = d\big(\vx^{(i)}, \vx^{(k)}\big), \qquad i, k = 1, \dots, N .
\]
Entry $d_{ik}$ answers one question: \emph{how dissimilar are stimuli $i$ and $k$, under
this measure, in this representation?} Row $i$ of $\mD$ is stimulus $i$'s dissimilarity to
everything.

\textbf{Example.} Euclidean, on the three images of \S4:
\[
\mD =
\begin{pmatrix}
0 & \sqrt{2} & \sqrt{18} \\
\sqrt{2} & 0 & \sqrt{20} \\
\sqrt{18} & \sqrt{20} & 0
\end{pmatrix}
\approx
\begin{pmatrix}
0 & 1.414 & 4.243 \\
1.414 & 0 & 4.472 \\
4.243 & 4.472 & 0
\end{pmatrix} .
\]

\subsection{Three properties every dissimilarity matrix has}

They follow from the three properties of a distance noted in \S4, and they are the first
thing to check when you have computed one:
\begin{itemize}
\item \textbf{Zero diagonal:} $d_{ii} = 0$. A stimulus is not dissimilar to itself.
\item \textbf{Symmetric:} $d_{ik} = d_{ki}$, so $\mD = \mD^{\top}$. The dissimilarity of
      $i$ to $k$ is the dissimilarity of $k$ to $i$.
\item \textbf{Non-negative:} $d_{ik} \geq 0$.
\end{itemize}
The same three hold for $\dcos$ and $\dcorr$; only the range of the off-diagonal entries
changes.

\subsection{Counting the pairs}

Because of the zero diagonal and the symmetry, an $N \times N$ dissimilarity matrix holds
only $N(N-1)/2$ distinct numbers, the entries above the diagonal (the \emph{upper triangle}).
For $N = 3$ that is $3$ numbers; for $N = 120$ it is $7{,}140$. When you compare two
dissimilarity matrices, those are the numbers you want: using all $N^2$ entries counts
every pair twice and adds $N$ zeros from the diagonal.

\np{np.allclose(D, D.T); np.allclose(np.diag(D), 0); (D >= 0).all().}

\section{Ranks and Spearman correlation}

Sometimes only the \emph{ordering} of a set of numbers is trustworthy, not their exact
values. Rating scales are like this: a ``7'' is above a ``6'' but not necessarily one
unit above. The \emph{rank} of an entry is its position when the entries are sorted from
smallest to largest.

\textbf{Example.} $\vx = (10, 20, 25)$ has ranks $(1, 2, 3)$; $\vy = (1, 5, 4)$ has ranks
$(1, 3, 2)$. Ties get the average of the positions they occupy: $(4, 4, 9)$ has ranks
$(1.5, 1.5, 3)$.

The \emph{Spearman correlation} $\rho$ is the Pearson correlation of the two rank vectors.
It measures whether two variables increase together, without assuming the relation is a
straight line. With no ties there is a shortcut:
\[
\rho = 1 - \frac{6\sum_j \delta_j^2}{n(n^2 - 1)}, \qquad
\delta_j = \text{rank}(x_j) - \text{rank}(y_j) .
\]
\textbf{Example.} Ranks $(1, 2, 3)$ and $(1, 3, 2)$ give $\bm{\delta} = (0, -1, 1)$,
$\sum \delta_j^2 = 2$, and $\rho = 1 - 6 \cdot 2 / (3 \cdot 8) = 0.5$. Computing Pearson
$r$ directly on the two rank vectors gives the same $0.5$: centre both to
$(-1, 0, 1)$ and $(-1, 1, 0)$, then $r = (1 + 0 + 0)/(\sqrt{2}\cdot\sqrt{2}) = 0.5$.

Two dissimilarity matrices are compared by taking the $N(N-1)/2$ upper-triangle entries of
each, in the same pair order, and computing $\rho$ between the two lists: a high $\rho$
says the pairs one representation calls far apart, the other does too, whatever the units.

\np{scipy.stats.spearmanr(a, b).statistic.}

\subsection*{Symbols used in this handout}
\begin{center}
\footnotesize
\setlength{\tabcolsep}{4pt}
\begin{tabular}{@{}ll@{\quad}ll@{}}
\toprule
$\vx$, $x_j$ & a vector; its $j$-th entry & $\norm{\vx}$ & Euclidean norm, $\sqrt{\vx \cdot \vx}$ \\
$D$, $N$ & entries per vector; number of stimuli & $\phi$, $\cos\phi$ & angle between vectors; cosine similarity \\
$\vx^{(i)}$, $x^{(i)}_j$ & response to stimulus $i$; entry $j$ & $\bar{x}$, $\tilde{\vx}$ & mean of $\vx$; $\vx$ centred by its own mean \\
$\mX$ & $N \times D$ response matrix & $r$, $\rho$ & Pearson; Spearman (Pearson on ranks) \\
$\vx \cdot \vy$ & dot product, $\sum_j x_j y_j$ & $\deuc$, $\dcos$, $\dcorr$ & $\norm{\vx - \vy}$, $1 - \cos\phi$, $1 - r$ \\
$\bm{A}^{\top}$ & transpose, rows and columns swapped & $\mD$, $d_{ik}$ & $N \times N$ dissimilarity matrix; entry $(i,k)$ \\
\bottomrule
\end{tabular}
\end{center}

\end{document}
